On a small MLP with symmetric label noise, SAM improved clean test accuracy over plain SGD by a large margin (\rhval{main/sam/digits_n0.2/test_acc/mean:3} vs \rhval{main/sgd/digits_n0.2/test_acc/mean:3} and \rhval{main/sam/digits_n0.4/test_acc/mean:3} vs \rhval{main/sgd/digits_n0.4/test_acc/mean:3} at 20\% and 40\% noise on digits), but weight decay alone recovered most of that gain, leaving a margin of \rhval{cmp/main/sgd+wd/digits_n0.2/test_acc/delta:4} at 20\% noise (inconclusive) and \rhval{cmp/main/sgd+wd/digits_n0.4/test_acc/delta:3} at 40\% (both tests below 0.05 uncorrected with $n=5$, one seed negative). SAM gave no benefit on noisy spirals, was sensitive to $\rho$, with a collapse to chance one grid step above the best value on noisy digits, and a random-direction perturbation of the same radius behaved like SGD. A fixed-radius loss-increase sharpness did not track the generalisation gap across optimisers (pooled Spearman \rhval{derived/corr-pooled-all/digits_n0.4/adv_loss/mean:2}--\rhval{derived/corr-pooled-all/digits_n0.4/rand_acc/mean:2}, negative within spirals). Within the limits of a five-seed, epoch-matched CPU study with untuned learning rate, the practical take-away is to compare SAM against a tuned weight decay and to tune $\rho$ on validation data; we make no claim about larger models.
